\documentclass[12pt]{article}
\usepackage[a4paper,margin=1in]{geometry}
\usepackage{amsmath, amssymb, booktabs, setspace}
\usepackage{enumitem}
\usepackage{graphicx}
\usepackage{siunitx}

\title{Practice Problem Set 4\\ \vspace{10pt} \large ECON 101 -- Macroeconomics}
\author{}
\vspace{-10pt}
\date{ \vspace{-30pt} Winter 2026}

\begin{document}
\maketitle

%%%%%%%%%%%%%%%%%%%%%%%%%%%%%%%%%%%%%%%%%%%%%%
\section*{Part I: Chapter 6 — Economic Growth and the Financial System}


\begin{enumerate}[label=\textbf{\arabic*.}]

\item Explain the difference between:
\begin{itemize}
    \item economic growth,
    \item business cycles.
\end{itemize}
Provide one real-world example of each.

\item Why is \textbf{labour productivity} crucial for long-run economic growth? 
Explain using:
\begin{itemize}
    \item capital deepening,
    \item technological progress.
\end{itemize}

\item Explain how \textbf{financial markets and financial intermediaries} contribute to economic growth. 
What would happen if financial markets were inefficient?

\end{enumerate}


\begin{enumerate}[label=\textbf{\arabic*.}, start=4]

\item \textbf{Loanable Funds Model}

Suppose:
\begin{itemize}
    \item National saving = \$800 billion,
    \item Investment demand = \$800 billion.
\end{itemize}

Now the government runs a deficit of \$200 billion.

\begin{enumerate}
    \item[(a)] What happens to national saving?
    \item[(b)] Show how the supply of loanable funds changes.
    \item[(c)] What happens to the equilibrium interest rate?
    \item[(d)] Explain crowding out.
\end{enumerate}
\newpage
\item \textbf{Growth Accounting}

Suppose:
\begin{itemize}
    \item Real GDP grows by 4\%,
    \item Labour input grows by 1\%.
\end{itemize}

\begin{enumerate}
    \item[(a)] Calculate labour productivity growth.
    \item[(b)] What factors could explain this growth?
\end{enumerate}

\end{enumerate}



\begin{enumerate}[label=\textbf{\arabic*.}, start=6]

\item Why do some developing countries fail to converge to the income levels of developed countries? 
Discuss at least three structural reasons.

\end{enumerate}

\hrule
\vspace{1em}

%%%%%%%%%%%%%%%%%%%%%%%%%%%%%%%%%%%%%%%%%%%%%%
\section*{Part II: Chapter 8 — Aggregate Expenditure and Output}


\begin{enumerate}[label=\textbf{\arabic*.}, start=7]

\item Define:
\begin{itemize}
    \item Marginal Propensity to Consume (MPC),
    \item Marginal Propensity to Save (MPS).
\end{itemize}
Explain their relationship.

\item Why does the \textbf{multiplier effect} amplify changes in spending? 
Explain intuitively.

\end{enumerate}


\begin{enumerate}[label=\textbf{\arabic*.}, start=9]

\item \textbf{Multiplier Calculation}

Suppose:
\begin{itemize}
    \item MPC = 0.8,
    \item Increase in government spending = \$50 billion.
\end{itemize}

\begin{enumerate}
    \item[(a)] Calculate the multiplier.
    \item[(b)] Calculate the change in equilibrium GDP.
\end{enumerate}

\item \textbf{Aggregate Expenditure Model}

Suppose:
\[
C = 200 + 0.75Y, \quad I = 150, \quad G = 250
\]

\begin{enumerate}
    \item[(a)] Write the aggregate expenditure (AE) equation.
    \item[(b)] Find equilibrium GDP.
    \item[(c)] Calculate the multiplier.
    \item[(d)] If \(G\) increases by 40, find the new equilibrium GDP.
\end{enumerate}

\end{enumerate}



\begin{enumerate}[label=\textbf{\arabic*.}, start=11]
\newpage

\item \textbf{Recessionary Gap}

\begin{itemize}
    \item Potential GDP = 3,000,
    \item Actual GDP = 2,600,
    \item MPC = 0.8.
\end{itemize}

\begin{enumerate}
    \item[(a)] Calculate the GDP gap.
    \item[(b)] Calculate the required increase in government spending to close the gap.
    \item[(c)] Explain the policy intuition.
\end{enumerate}

\end{enumerate}

\hrule
\vspace{1em}

%%%%%%%%%%%%%%%%%%%%%%%%%%%%%%%%%%%%%%%%%%%%%%
\section*{Part III: Chapter 9 — Aggregate Demand and Aggregate Supply}



\begin{enumerate}[label=\textbf{\arabic*.}, start=12]

\item Explain why the aggregate demand (AD) curve is downward sloping using:
\begin{itemize}
    \item wealth effect,
    \item interest rate effect,
    \item international trade effect.
\end{itemize}

\item Distinguish between:
\begin{itemize}
    \item short-run aggregate supply (SRAS),
    \item long-run aggregate supply (LRAS).
\end{itemize}

\item What causes shifts in:
\begin{itemize}
    \item AD,
    \item SRAS,
    \item LRAS?
\end{itemize}
Give one example for each.

\end{enumerate}


\begin{enumerate}[label=\textbf{\arabic*.}, start=15]

\item \textbf{Demand Shock}

Suppose consumer confidence falls.

\begin{enumerate}
    \item[(a)] What happens to AD?
    \item[(b)] What are the short-run effects on:
    \begin{itemize}
        \item output,
        \item price level,
        \item unemployment?
    \end{itemize}
\end{enumerate}
\newpage
\item \textbf{Supply Shock}

Suppose oil prices increase significantly.

\begin{enumerate}
    \item[(a)] What happens to SRAS?
    \item[(b)] What happens to inflation and output?
    \item[(c)] What type of inflation is this?
\end{enumerate}

\end{enumerate}



\begin{enumerate}[label=\textbf{\arabic*.}, start=17]

\item Suppose:
\begin{itemize}
    \item Potential GDP = 2,500,
    \item Actual GDP = 2,700,
    \item Inflation is rising.
\end{itemize}

\begin{enumerate}
    \item[(a)] Is there an inflationary or recessionary gap?
    \item[(b)] What policy should be used?
    \item[(c)] Explain the adjustment process back to equilibrium.
\end{enumerate}

\end{enumerate}



\begin{enumerate}[label=\textbf{\arabic*.}, start=18]

\item The economy experiences:
\begin{itemize}
    \item a decline in investment,
    \item rising interest rates,
    \item falling GDP.
\end{itemize}

\begin{enumerate}
    \item[(a)] Explain using:
    \begin{itemize}
        \item AE model,
        \item AD--AS model.
    \end{itemize}
    \item[(b)] What role does the financial system play in this downturn?
    \item[(c)] Propose one fiscal policy and one monetary policy. 
    Explain how each affects GDP.
\end{enumerate}

\end{enumerate}



\begin{enumerate}[label=\textbf{\arabic*.}, start=19]

\item Suppose:
\begin{itemize}
    \item Government spending increases by \$100,
    \item MPC = 0.75,
    \item Interest rates rise due to deficit financing.
\end{itemize}

Explain why the actual increase in GDP may be smaller than predicted by the multiplier.

\end{enumerate}

\hrule
\vspace{1em}

\begin{center}
\textit{End of Problem Set}
\end{center}

\end{document}